\subsubsection{Subproceso 2: Selección y Registro de Cursos}




El Subproceso 2 tiene como propósito permitir al estudiante habilitado conformar y registrar su propuesta académica semestral en la Intranet DIRCE, seleccionando las asignaturas, secciones y horarios de su Plan de Estudios dentro de los límites de creditaje autorizados por su promedio ponderado y durante la franja horaria de turno asignada en el Subproceso 1.

Este subproceso constituye el núcleo operativo de la matrícula regular: es el único momento en que el estudiante ejerce su autonomía para construir su carga académica del periodo. Actúa como interfaz entre la oferta de la Facultad y la demanda individual, con la plataforma SIGA-ORCE como árbitro automático de concurrencia y validación en tiempo real. Su ejecución exitosa produce la Pre-Ficha de Matrícula Registrada, insumo indispensable para los subprocesos posteriores de auditoría y confirmación oficial.



